\clearpage
\fancyhead[L]{\small MAGNETIC EFFECT AT A DISTANCE.}
\SourceLeaf{433}\pdfbookmark[0]{Instructions for determining the magnetic effect exerted by a bar magnet at a distance.}{fernwirkung-title}
\begin{center}
{\Large\bfseries INSTRUCTIONS FOR DETERMINING\par\medskip THE MAGNETIC EFFECT EXERTED\par\medskip BY A BAR MAGNET\par\medskip AT A DISTANCE.}
\par\bigskip
Resultate aus den Beobachtungen des magnetischen Vereins. 1840. II.
\end{center}
\renewcommand{\Article}[1]{\par\medskip\Needspace{4\baselineskip}\pdfbookmark[1]{#1.}{fernwirkung-art-#1}\centerline{\textbf{#1.}}\nobreak\smallskip}%
When several magnetic instruments are set up at the same time, their effects on one another must be taken into account. Placing the instruments so far apart that these effects can be assumed without further examination to be entirely imperceptible is a way around the problem that cannot be used everywhere and in any case entails giving up some other advantages and conveniences. If, however, we can calculate precisely the effects that one instrument exerts at the site of another, and thus separate them from the observations made with the second instrument, we are entirely free to give every other consideration its due when choosing the sites. The derivation of the formulas for these calculations therefore seems well worth including here.

\Article{1}
Let the position of the point at which the effect of a bar magnet is to be calculated be given by three rectangular coordinates, $x,y,z$, with the origin here placed at the centre of the bar magnet itself. To fix ideas, we assume that the first two coordinate axes are horizontal, the first in the true meridian, and that the third is
\SourceLeaf{434}therefore vertical. We count $x$ positive towards the south, $y$ towards the west and $z$ upwards. We also set
\[
\BareRoot{(xx+yy+zz)}=r,\quad x=r\cos f\cos g,\quad y=r\cos f\sin g,\quad z=r\sin f
\]
so that $g$ denotes the azimuth of the straight line drawn from the centre of the bar to the point in question and $f$ its inclination to the horizontal plane.

Further, let $M$ denote the absolute magnetic moment of the bar magnet, $F$ the inclination of its magnetic axis, taken as positive when the north pole is higher, and $G$ the azimuth of that axis. For brevity we write
\[
\cos F\cos G=A,\qquad\cos F\sin G=B,\qquad\sin F=C
\]
so the magnetic moments of the bar magnet with respect to the three coordinate axes are $MA,MB,MC$, respectively.

We resolve the magnetic force exerted by the bar magnet at the point $x,y,z$ into the component forces $\xi,\eta,\zeta$ parallel to the three coordinate axes.

Let $U$ denote the total intensity of the purely geomagnetic force at this location and $Z$ its vertical component. We resolve the horizontal component $T$ into the component forces $X$ and $Y$ parallel to the first two coordinate axes. All these forces $\xi,\eta,\zeta,T,U,X,Y,Z$ are quantities of the same kind.

Finally, let $i$ be the magnetic inclination and $D$ the declination. To follow ordinary usage, we count $D$ from north towards west and regard $i$ as positive when the north pole of the magnetic needle dips downwards. Thus we have
\[
\begin{aligned}
X&=-T\cos D=-U\cos i\cos D\\
Y&=T\sin D=U\cos i\sin D\\
Z&=-T\operatorname{tang}i=-U\sin i
\end{aligned}
\]

\Article{2}
The effect of the bar magnet at the point $x,y,z$ consists in small changes in the quantities that determine the geomagnetic force. Because these changes are small, we can safely treat them according to the rules of differential calcu%
\SourceLeaf{435}lus, and we shall denote them by prefixing the symbol $\mathrm{d}$. Since
\[
\mathrm{d}X=\xi,\qquad\mathrm{d}Y=\eta,\qquad\mathrm{d}Z=\zeta
\]
we have
\[
\begin{aligned}
\xi&=T\sin D.\mathrm{d}D-\cos D.\mathrm{d}T\\
\eta&=T\cos D.\mathrm{d}D+\sin D.\mathrm{d}T\\
\zeta&=-T\sec i^2\mathrm{d}i-\operatorname{tg}i.\mathrm{d}T
\end{aligned}
\]
from which we obtain
\[
\begin{aligned}
\mathrm{d}D&=\frac{\sin D}{T}.\xi+\frac{\cos D}{T}.\eta\\
\mathrm{d}T&=-\cos D.\xi+\sin D.\eta\\
\mathrm{d}i&=-\frac{\cos i^2}{T}.\zeta-\frac{\sin2i}{2T}.\mathrm{d}T
\end{aligned}
\]
Finally, we have
\[
\mathrm{d}U=\cos i.\mathrm{d}T-\sin i.\zeta
\]
or
\[
\frac{\mathrm{d}U}{U}=\frac{\cos i^2}{T}.\mathrm{d}T-\frac{\sin2i}{2T}.\zeta=\frac{\mathrm{d}T}{T}+\operatorname{tang}i.\mathrm{d}i
\]

\Article{3}
The potential at the point $x,y,z$ of the magnetic fluids contained in the bar magnet can be expanded in a series in powers of $\frac1r$. For our purpose we need to retain only the leading term, which is of order $\frac{1}{rr}$. If we denote this potential by $V$, it is easy to see that under this restriction we have
\[
V=\frac{M(Ax+By+Cz)}{r^3}
\]
As is well known, $\xi,\eta,\zeta$ are obtained from the partial derivatives of $V$ with respect to $x,y,z$. We have
\[
\xi=-\frac{\mathrm{d}V}{\mathrm{d}x},\qquad\eta=-\frac{\mathrm{d}V}{\mathrm{d}y},\qquad\zeta=-\frac{\mathrm{d}V}{\mathrm{d}z}
\]
\SourceLeaf{436}and consequently, if for brevity we set
\[
\frac{Ax+By+Cz}{r}=k
\]
and note that the partial derivatives $\frac{\mathrm{d}r}{\mathrm{d}x},\frac{\mathrm{d}r}{\mathrm{d}y},\frac{\mathrm{d}r}{\mathrm{d}z}$ are respectively $=\frac{x}{r},\frac{y}{r},\frac{z}{r}$, we obtain
\[
\begin{aligned}
\xi&=\frac{3Mkx}{r^4}-\frac{MA}{r^3}\\
\eta&=\frac{3Mky}{r^4}-\frac{MB}{r^3}\\
\zeta&=\frac{3Mkz}{r^4}-\frac{MC}{r^3}
\end{aligned}
\]
Substituting the values of $x,y,z,A,B,C$, we thus obtain
\[
\begin{aligned}
k&=\cos f\cos F\cos(G-g)+\sin f\sin F\\
\mathrm{d}D&=\frac{M}{Tr^3}(3k\cos f\sin(D+g)-\cos F\sin(D+G))\\
\frac{\mathrm{d}T}{T}&=\frac{M}{Tr^3}(-3k\cos f\cos(D+g)+\cos F\cos(D+G))\\
\mathrm{d}i&=-\tfrac12\sin2i.\frac{\mathrm{d}T}{T}-\frac{M}{Tr^3}.\cos i^2(3k\sin f-\sin F)\\
\frac{\mathrm{d}U}{U}&=\cos i^2.\frac{\mathrm{d}T}{T}-\frac{M}{2Tr^3}.\sin2i(3k\sin f-\sin F)
\end{aligned}
\]
These formulas give the complete solution to our problem. It is clear without our pointing it out that $\mathrm{d}D$ and $\mathrm{d}i$ are expressed here in units of the radius, so the factor $206265''$ must still be applied to the values to obtain those changes in units of arc. The value of $\frac{M}{T}$ will have to be determined by experiments using the method given in \textit{Intensitas vis magneticae terrestris}.

\Article{4}
In practice, the most common cases are those in which special circumstances considerably simplify our general formulas. The following two in particular deserve mention.

I. If the bar magnet is vertical, so that $F=\pm90^0$, the general formulas take the following form:
\SourceLeaf{437}
\[
\begin{aligned}
k&=\pm\sin f\\
\mathrm{d}D&=\pm\frac{3M}{2Tr^3}\sin2f\sin(D+g)\\
\frac{\mathrm{d}T}{T}&=\mp\frac{3M}{2Tr^3}\sin2f\sin(D+g)\\
\mathrm{d}i&=-\tfrac12\sin2i.\frac{\mathrm{d}T}{T}\mp\frac{M}{Tr^3}\cos i^2(3\sin f^2-1)\\
\frac{\mathrm{d}U}{U}&=\cos i^2.\frac{\mathrm{d}T}{T}\mp\frac{M}{2Tr^3}\sin2i(3\sin f^2-1)
\end{aligned}
\]
If the point $x,y,z$ is also at the same height as the centre of the bar magnet, then $z=0,f=0$, and consequently
\[
\begin{aligned}
\mathrm{d}D&=0\\
\mathrm{d}T&=0\\
\mathrm{d}i&=\pm\frac{M}{Tr^3}\cos i^2\\
\frac{\mathrm{d}U}{U}&=\pm\frac{M}{2Tr^3}\sin2i
\end{aligned}
\]
It follows that observations with a unifilar or bifilar magnetometer are entirely undisturbed by a second bar magnet in the same room, as long as that bar is kept vertical and its centre at the same height as the magnetometer bar.

II. If the bar magnet is horizontal, or $F=0$, our formulas become:
\[
\begin{aligned}
k&=\cos f\cos(G-g)\\
\mathrm{d}D&=\frac{M}{Tr^3}(3\cos f^2\cos(G-g)\sin(D+g)-\sin(D+G))\\
\frac{\mathrm{d}T}{T}&=\frac{M}{Tr^3}(-3\cos f^2\cos(G-g)\cos(D+g)+\cos(D+G))\\
\mathrm{d}i&=-\tfrac12\sin2i.\frac{\mathrm{d}T}{T}-\frac{3M}{2Tr^3}\cos i^2\sin2f\cos(G-g)\\
\frac{\mathrm{d}U}{U}&=\cos i^2.\frac{\mathrm{d}T}{T}-\frac{3M}{4Tr^3}\sin2i\sin2f\cos(G-g)
\end{aligned}
\]
If the bar magnet also lies in the magnetic meridian (so that $G=180^0-D$ in the natural position) or perpendicular to it (so that $G=90^0-D$ or $=270^0-D$, according as the north pole is on the west or east side), the formulas for $\mathrm{d}D$ and $\mathrm{d}T$ evidently admit still further simpli%
\SourceLeaf{438}fication. These cases occur when the bar whose effect at a distance is sought forms part of a unifilar magnetometer or of a bifilar magnetometer in the transverse position.

\Article{5}
If we want to compare the effects of a bar magnet in different horizontal positions, we can easily put each of the formulas given in the preceding article II into a form suited to this purpose. If, for example, we determine two quantities $p,P$ by the equations
\[
\begin{aligned}
p\cos P&=(3\cos f^2-1)\sin(D+g)\\
p\sin P&=\cos(D+g)
\end{aligned}
\]
the formula for $\mathrm{d}D$ becomes
\[
\mathrm{d}D=\frac{Mp}{Tr^3}\cos(G-g+P)
\]
It follows that $\mathrm{d}D$, for $G=g-P$ or for $G=180^0+g-P$, reaches its greatest value $\frac{Mp}{Tr^3}$ with a positive or negative sign, whereas it vanishes for $G=90^0+g-P$ and for $G=270^0+g-P$. In the same way, if we set
\[
\begin{aligned}
q\cos Q&=(3\cos f^2-1)\cos(D+g)\\
q\sin Q&=\sin(D+g)
\end{aligned}
\]
we obtain
\[
\frac{\mathrm{d}T}{T}=-\frac{Mq}{Tr^3}\cos(G-g-Q)
\]
from which we can similarly determine its maximum value and when it vanishes.

\Article{6}
The arrangements here offer several opportunities to apply these instructions, both in determining the mutual effects of the bars of the unifilar and bifilar magnetometers on one another and in determining the effects of both bars at a third site, where magnetic observations are made outdoors from time to time with other instruments on a fixed stone pier. The absolute coordinates of these three sites, expressed in metres and referred to the centre
\SourceLeaf{439}of the axis of \textsc{Reichenbach}’s meridian circle and, for the third coordinate, to the floor of the observatory, are as follows:

(I) Centre of the twenty-five-pound bar magnet of the bifilar magnetometer, for which, with the metre taken as the unit of length, $\frac{M}{T}=2.63318$,
\[
x=-3.391,\qquad y=+6.708,\qquad z=+0.661
\]
(II) Centre of the four-pound bar magnet of the unifilar magnetometer, for which $\frac{M}{T}=0.48592$
\[
x=-23.618,\qquad y=+69.206,\qquad z=-2.235
\]
(III) Centre of the stone pier and, as regards height, the position occupied by the centre of the needle of a \textsc{Robinson} dip circle,
\[
x=-21.546,\qquad y=+56.979,\qquad z=-1.665
\]
Only the final results of a calculation for four cases are given here. The values used for $D$ and $i$ are $18^0 11'54''$ and $67^0 36'$. For the present purpose the variation of these quantities, as well as of the values of $\frac{M}{T}$ for the two bar magnets, need not be considered.

(1) and (4) Effects of the bar magnet of the unifilar magnetometer, the former at site (III) and the latter at site (I).

(2) and (3) Effects of the bar magnet of the bifilar magnetometer at sites (III) and (II), assuming that bar to be in the transverse position with its north pole to the west.
\[
\begin{array}{c|r|r|r|r}
&\mathrm{d}D&\mathrm{d}T&\mathrm{d}i&\mathrm{d}U\\\hline
(1)&+64''72&-0.0000884\,T&+6''91&-0.0000071\,U\\
(2)&+3.04&+0.0000250\,T&-1.76&+0.0000043\,U\\
(3)&+1.82&+0.0000132\,T&-0.93&+0.0000023\,U\\
(4)&+0.50&+0.0000001\,T&-0.00&+0.0000001\,U
\end{array}
\]
The numbers for (2) and (3) merely change sign when the bar in the bifilar magnetometer is in the transverse position with its north pole to the east. The total disturbance at site III from both instruments is therefore
\SourceLeaf{440}
\[
\begin{array}{c|r|r}
&\multicolumn{2}{c}{\text{North pole in the bifilar magnetometer.}}\\
&\text{West}&\text{East}\\\hline
\mathrm{d}D&+67''76&+61''68\\
\mathrm{d}T&-0.0000634\,T&-0.0001134\,T\\
\mathrm{d}i&+5''15&+8''67\\
\mathrm{d}U&-0.0000028\,U&-0.0000114\,U
\end{array}
\]

\Article{7}
Finally, we shall show the connection between the formulas given in Article 2 for the effect of a bar magnet at a distance and a simple construction already mentioned in volume 2 of the \textit{Resultate} [1837. II.]. The reader can easily draw a figure from the instructions below or consult it in the place cited.

Let $A$ be the centre of the bar magnet, $n$ any other point on its magnetic axis through $A$ on the side of the magnetic north pole, $s$ a similar point on the side of the south pole, and $C$ the point at which we are to determine the magnetic effect of the bar on the unit of northern magnetic fluid imagined to be concentrated there. According to Article 2, the component forces $\xi,\eta,\zeta$ are expressed by the formulas
\[
\begin{aligned}
\xi&=\frac{3Mkx}{r^4}-\frac{MA}{r^3}\\
\eta&=\frac{3Mky}{r^4}-\frac{MB}{r^3}\\
\zeta&=\frac{3Mkz}{r^4}-\frac{MC}{r^3}
\end{aligned}
\]
where, as is easy to see, $k$ equals the cosine of the angle between $An$ and $AC$. The first parts of $\xi,\eta,\zeta$ evidently combine into a single force $\frac{3Mk}{r^3}$, which acts repulsively in the direction of the straight line $AC$ when $k$ is positive and attractively, or in the opposite direction $CA$, when $k$ is negative. Likewise the second parts of $\xi,\eta,\zeta$ combine into a single force $\frac{M}{r^3}$ whose direction is always parallel to $ns$. In the special case in which $AC$ makes a right angle with the magnetic axis, so that $k=0$, the first force vanishes and the second alone gives the
\SourceLeaf{441}entire effect. In every other case, in the plane containing $n,A,s,C$, let $CB$ be a perpendicular to $CA$, $B$ its intersection with the magnetic axis, and [$D$ the point on the straight line $AB$ for which] $AD=\tfrac13AB$. In the case of the figure in volume 2 of the \textit{Resultate}, where $AC$ makes an obtuse angle with $An$, so that $D$ and $B$ lie on the side of the south pole, the two forces given above are evidently proportional to the straight lines $CA$ and $AD$. The direction of the first coincides with $CA$, and the direction of the other is parallel to $AD$. The direction of their resultant will therefore be $CD$ and its magnitude $=\frac{CD}{AD}.\frac{M}{r^3}$. In the other case, which is absent from the figure in the place cited, $AC$ makes an acute angle with $An$, so that $B$ and $D$ lie on the side of the north pole. The same result holds with the sole difference that the direction of the angle of the bar magnet on an element of northern fluid is not expressed by $CD$, but by $DC$. This must therefore still be added in the place cited for completeness.
